% 教学构造：柱高只展示偏差、方差与不可约噪声的相对变化。
\begin{tikzpicture}[foundation picture,x=10mm,y=10mm,
  font=\figurefont\footnotesize]
  \draw[foundation axis] (0,0)--(8.2,0) node[right] {模型复杂度};
  \draw[foundation axis] (0,0)--(0,6.4) node[above] {期望预测误差};
  \foreach \i/\bias/\variance in {
    1/4.5/.3,2/3.5/.5,3/2.5/.8,4/1.7/1.2,
    5/1.0/1.9,6/.6/2.8,7/.4/4.0}{
    \fill[FigureBlue] ({\i-.28},0) rectangle ({\i+.28},\bias);
    \fill[FigureRed] ({\i-.28},\bias)
      rectangle ({\i+.28},{\bias+\variance});
    \fill[FigureYellow] ({\i-.28},{\bias+\variance})
      rectangle ({\i+.28},{\bias+\variance+.8});
    \draw[FigureStroke,line width=.5pt] ({\i-.28},0)
      rectangle ({\i+.28},{\bias+\variance+.8});
  }
  \node[fill=none,align=left,inner sep=2pt] at (6.7,6.0)
    {\textcolor{FigureBlue}{\rule{3mm}{2mm}} 偏差平方\\
     \textcolor{FigureRed}{\rule{3mm}{2mm}} 方差\\
     \textcolor{FigureYellow}{\rule{3mm}{2mm}} 不可约噪声};
  \draw[FigureStroke,dashed,line width=.8pt] (0.4,.8)--(7.6,.8);
  \node[anchor=north] at (1,-.25) {低};
  \node[anchor=north] at (7,-.25) {高};
\end{tikzpicture}
